\documentclass[UTF8,zihao=-4]{ctexart}
\usepackage[a4paper,margin=2.2cm]{geometry}
\usepackage{xcolor}
\usepackage{hyperref}
\usepackage{enumitem}
\usepackage{fontspec}
\usepackage{amsmath}
\IfFontExistsTF{Microsoft YaHei}{\setCJKmainfont{Microsoft YaHei}}{\setCJKmainfont{SimSun}[BoldFont=SimHei]}
\IfFontExistsTF{Microsoft YaHei UI}{\setCJKsansfont{Microsoft YaHei UI}}{}
\IfFontExistsTF{Consolas}{\setmonofont{Consolas}}{}
\definecolor{linkblue}{RGB}{30,80,145}
\hypersetup{colorlinks=true,linkcolor=linkblue,urlcolor=linkblue,pdfborder={0 0 0}}
\setlength{\parindent}{2em}
\setlength{\parskip}{0.25em}
\linespread{1.16}
\setlist{itemsep=0.2em,topsep=0.35em,leftmargin=2.5em}
\pagestyle{plain}

\title{05\_堆的定义数组表示与操作}

\author{}

\date{}

\begin{document}

\maketitle

\section*{05 堆的定义、数组表示与操作}

\subsection*{题目原文}

了解堆的定义、如何用数组来表示堆和实现堆上的操作。

\subsection*{考查类型}

概念题和算法设计题。

\subsection*{堆的定义}

堆是一棵满足堆序性质的完全二叉树。

最大堆：

\begin{quote}
\small\ttfamily\raggedright
\relax 每个结点的值\ ≥\ 它的子结点的值\\
\end{quote}

最小堆：

\begin{quote}
\small\ttfamily\raggedright
\relax 每个结点的值\ ≤\ 它的子结点的值\\
\end{quote}

完全二叉树保证堆可以紧凑地存入数组，不需要显式保存指针。

\subsection*{数组表示}

若数组从下标 \texttt{1} 开始：

\[parent(i) = \lfloor i / 2 \rfloor\]
\[left(i) = 2i\]
\[right(i) = 2i + 1\]

若数组从下标 \texttt{0} 开始：

\[parent(i) = \lfloor (i - 1) / 2 \rfloor\]
\[left(i) = 2i + 1\]
\[right(i) = 2i + 2\]

考试中要先说明使用哪一种下标体系，再写公式。

\subsection*{主要操作}

\subsubsection*{向上调整}

插入新元素时，把元素放在数组末尾。若它破坏堆序性质，就不断与父结点交换，直到堆序恢复。

时间复杂度：

\begin{quote}
\small\ttfamily\raggedright
\relax O(log\ n)\\
\end{quote}

\subsubsection*{向下调整}

删除堆顶时，用最后一个元素补到堆顶。若它破坏堆序性质，就不断与更合适的子结点交换，直到堆序恢复。

最大堆中应与较大的子结点交换。最小堆中应与较小的子结点交换。

时间复杂度：

\begin{quote}
\small\ttfamily\raggedright
\relax O(log\ n)\\
\end{quote}

\subsubsection*{建堆}

自底向上对所有非叶子结点做向下调整。

最后一个非叶子结点的位置为：

\[\lfloor n / 2 \rfloor\]

这里采用 \texttt{1} 开始的数组下标。

总时间复杂度：

\begin{quote}
\small\ttfamily\raggedright
\relax O(n)\\
\end{quote}

\subsection*{典型伪代码}

最大堆的向下调整：

\begin{quote}
\small\ttfamily\raggedright
\relax MaxHeapify(A,\ i,\ heapSize):\\
\end{quote}
\[l = left(i)\]
\[r = right(i)\]
\[largest = i\]
\begin{quote}
\small\ttfamily\raggedright
\relax \ \ \ \ if\ l\ ≤\ heapSize\ and\ A[l]\ >\ A[largest]:\\
\end{quote}
\[largest = l\]
\begin{quote}
\small\ttfamily\raggedright
\relax \ \ \ \ if\ r\ ≤\ heapSize\ and\ A[r]\ >\ A[largest]:\\
\end{quote}
\[largest = r\]
\begin{quote}
\small\ttfamily\raggedright
\relax \ \ \ \ if\ largest\ ≠\ i:\\
\relax \ \ \ \ \ \ \ \ swap\ A[i],\ A[largest]\\
\relax \ \ \ \ \ \ \ \ MaxHeapify(A,\ largest,\ heapSize)\\
\end{quote}

\subsection*{常见应用}

优先队列可以用堆实现。

堆排序可以用最大堆实现：先建最大堆，再反复取出堆顶最大元素。

\subsection*{例题}

\subsubsection*{例题 1}

题目：判断数组是否表示一个最大堆。数组从下标 \texttt{1} 开始：

\begin{quote}
\small\ttfamily\raggedright
\relax A\ =\ [16,\ 14,\ 10,\ 8,\ 7,\ 9,\ 3,\ 2,\ 4,\ 1]\\
\end{quote}

解答：

逐个检查非叶子结点：

\[16 \ge  14, 10\]
\[14 \ge  8, 7\]
\[10 \ge  9, 3\]
\[8 \ge  2, 4\]
\[7 \ge  1\]

每个父结点都不小于子结点，所以该数组表示一个最大堆。

\subsubsection*{例题 2}

题目：向上面的最大堆插入 \texttt{15}，写出调整后的数组。

解答：

先把 \texttt{15} 放到数组末尾：

\begin{quote}
\small\ttfamily\raggedright
\relax [16,\ 14,\ 10,\ 8,\ 7,\ 9,\ 3,\ 2,\ 4,\ 1,\ 15]\\
\end{quote}

\texttt{15} 的父结点是 \texttt{7}，交换：

\begin{quote}
\small\ttfamily\raggedright
\relax [16,\ 14,\ 10,\ 8,\ 15,\ 9,\ 3,\ 2,\ 4,\ 1,\ 7]\\
\end{quote}

\texttt{15} 的新父结点是 \texttt{14}，继续交换：

\begin{quote}
\small\ttfamily\raggedright
\relax [16,\ 15,\ 10,\ 8,\ 14,\ 9,\ 3,\ 2,\ 4,\ 1,\ 7]\\
\end{quote}

此时 $15 \le 16$，调整结束。

\subsubsection*{例题 3}

题目：从原最大堆删除堆顶 \texttt{16}，写出调整后的数组。

解答：

用末尾元素 \texttt{1} 补到堆顶，再向下调整：

\begin{quote}
\small\ttfamily\raggedright
\relax [1,\ 14,\ 10,\ 8,\ 7,\ 9,\ 3,\ 2,\ 4]\\
\relax [14,\ 1,\ 10,\ 8,\ 7,\ 9,\ 3,\ 2,\ 4]\\
\relax [14,\ 8,\ 10,\ 1,\ 7,\ 9,\ 3,\ 2,\ 4]\\
\relax [14,\ 8,\ 10,\ 4,\ 7,\ 9,\ 3,\ 2,\ 1]\\
\end{quote}

最终最大堆为：

\begin{quote}
\small\ttfamily\raggedright
\relax [14,\ 8,\ 10,\ 4,\ 7,\ 9,\ 3,\ 2,\ 1]\\
\end{quote}

\subsection*{常见误区}

堆不是二叉搜索树。堆只保证父子之间的大小关系，不保证左子树所有元素都小于右子树所有元素。

建堆不是 $O(n \log n)$ 的紧分析结果。逐个插入建堆是 $O(n \log n)$，自底向上建堆是 $O(n)$。

\end{document}
